\section{Design}
\label{sec:design}

We now describe the design of \sysname{}, a GPU-native I/O runtime
that executes storage operations directly on the GPU.
\sysname{} comprises three layers: (1)~a \emph{task execution engine}
that manages asynchronous, warp-level coroutine execution on the GPU
(\S\ref{sec:exec-engine}); (2)~a \emph{dynamic work orchestrator} that
maps tasks of varying parallelism to GPU warps and adapts resource
allocation to load (\S\ref{sec:orchestrator}); and (3)~\emph{AggroFS},
a GPU-resident blob-store filesystem that performs metadata management
and storage tiering entirely on the GPU (\S\ref{sec:aggrofs}).

Figure~\ref{fig:arch-overview} depicts the end-to-end architecture,
illustrating the GPU$\rightarrow$GPU, GPU$\rightarrow$CPU, and
CPU$\rightarrow$GPU communication paths.

\begin{figure}[t]
  \centering
  \includesvg[width=\columnwidth]{fig_arch_overview}
  \caption{%
    Architecture overview of \sysname{}.  The GPU runtime hosts
    warp-level workers that execute module tasks as coroutines.
    Clients submit tasks through per-warp queues mapped in device
    memory; the IPC manager routes responses back via shared memory,
    IPC sockets, or TCP\@.  CPU-side modules (e.g., NVMe drivers) are
    reached through pinned-host queues.}
  \label{fig:arch-overview}
\end{figure}

%%%%============================================================
\subsection{Asynchronous GPU Task Execution Engine}
\label{sec:exec-engine}
%%%%============================================================

\paragraph{Runtime startup.}
The CPU runtime is launched first as a privileged process.
It initializes NVMe submission and completion queues and maps them
into the GPU address space using BAR-mapped MMIO regions, enabling
the GPU to issue storage commands without CPU involvement.
The CPU runtime then spawns the GPU runtime as a persistent CUDA
kernel.  During initialization, the \texttt{IpcManager} allocates a
configurable number of \emph{per-warp task queues} in device memory---one
set for CPU$\rightarrow$GPU traffic and one for GPU$\rightarrow$GPU
traffic---along with a \emph{suspended-task queue} per warp for
coroutines that are waiting on I/O\@.
The number of warp queues determines the maximum concurrency of both
runtimes: each warp constitutes an independent, lock-free scheduling
domain.

\paragraph{Client connection.}
Once both runtimes are running, clients connect through the IPC
manager.  A client first allocates a region of \emph{device memory}
(for GPU-resident clients) or \emph{pinned host memory} (for CPU
clients) to hold communication primitives: task descriptors and
\texttt{FutureShm} completion tokens.
The IPC manager supports three transport modes---shared memory, Unix
domain sockets, and TCP---chosen at connection time via an environment
variable.  Regardless of transport, the client receives a pool of
pre-allocated \texttt{FutureShm} objects backed by shared or device
memory.  These objects carry all the metadata needed to route a
response back to the caller: the originating transport, a client
process identifier, and a flexible-array \texttt{copy\_space} for
small return values.

\paragraph{An efficient communication protocol.}
Task submission follows a zero-copy, poll-based protocol.
A client constructs a task descriptor in its private memory region,
populates a \texttt{PoolQuery} that specifies the target module,
routing mode, and degree of parallelism, and pairs the task with a
\texttt{FutureShm} token drawn from the pre-allocated pool.
Input data is registered as \emph{bulk descriptors} referencing
contiguous device- or host-memory regions; no copies occur at
submission time.  The client then enqueues the task pointer into one
of the per-warp queues using an atomic compare-and-swap, after which
the submitting thread is free to continue.

On the runtime side, lane~0 of each warp continuously polls its
queues.  When a task arrives, the warp instantiates a coroutine (see
below), executes the target module's entry point, and writes any
return values into the \texttt{FutureShm}'s \texttt{copy\_space}.
For shared-memory clients, the runtime sets the
\texttt{FUTURE\_COMPLETE} flag in the token and the client observes
completion by polling the flag.  For socket- and TCP-based clients,
the runtime additionally pushes a notification through the
corresponding transport.  Because the entire protocol is lock-free on
the fast path---using only atomic flags and single-producer
queues---it adds negligible overhead to the critical path of an I/O
request.

\paragraph{Warp-level coroutines.}
After a task reaches a warp, the worker wraps the module's entry
point in a \emph{warp-level coroutine}: a C++20 coroutine that
executes collectively across all 32 lanes of a single warp.
Standard NVCC does not support C++20 coroutines in device code;
we augment the Clang/LLVM CUDA front-end to lower
\texttt{co\_await} and \texttt{co\_yield} into device-compatible
suspend and resume points.

Each coroutine frame is bump-allocated from a per-warp stack
region in device memory; only when the stack is exhausted does the
allocator fall back to a shared \texttt{BuddyAllocator}.
A \texttt{RunContext} object, allocated once per active coroutine,
holds 32 \texttt{coroutine\_handle}s (one per lane), a yield flag,
and a pointer to the \texttt{FutureShm} being awaited.
When a coroutine issues a sub-task (e.g., an NVMe read), it
\texttt{co\_await}s the sub-task's \texttt{Future}, which suspends
all 32 lanes simultaneously and returns control to lane~0's
scheduler loop.  The scheduler then picks the next runnable
coroutine from the warp's suspended queue or polls for a new
incoming task.  When the awaited I/O completes---signaled by the
sub-task's \texttt{FutureShm} flag---the scheduler resumes the
suspended coroutine.

This design yields two key benefits.
First, it hides I/O latency: while one coroutine waits on a
device read, the warp executes another coroutine, keeping the
SIMT lanes productive.
Second, it preserves warp-coherent execution: because all 32
lanes enter and exit suspend points together, there is no
divergence penalty, and the programming model remains
single-instruction.

\paragraph{Module specification.}
A \emph{module} (called a \emph{ChiMod} internally) is a dynamically
loadable code object that implements a named service---for example,
\texttt{EXT4} for a POSIX filesystem or \texttt{BDEV} for a block
device.  Each module instance is identified by three attributes: a
\emph{library name} (e.g., \texttt{bdev}), a \emph{semantic instance
name} (e.g., a mount point or device path), and a
\emph{numeric instance ID} used for fast routing.

On the CPU, modules are represented as C++ objects with virtual
dispatch.
On the GPU, virtual function calls incur an indirection through
global memory that serializes warp execution; we therefore replace
the vtable with an \emph{inline function-pointer table} embedded
directly in the container object.  This reduces each module
invocation to a single global-memory load followed by a direct
call, eliminating the pointer-chase overhead of virtual dispatch.

Modules expose a small set of callbacks that the runtime invokes at
well-defined points: \texttt{Create} (initialize a new instance),
\texttt{Destroy} (tear down), \texttt{Run} (execute a task method),
and \texttt{Update} (refresh GPU-side state after a code or
configuration change).  The \texttt{Run} callback dispatches on a
method identifier embedded in the task, enabling a single module to
expose an arbitrarily rich interface while paying only one
indirect-call overhead.

\paragraph{Task specification.}
A \emph{task} is the unit of work submitted to a module.
Its base class carries the parameters needed to route and scale
execution:
a \texttt{pool\_id} identifying the target module instance,
a \texttt{method} selector for the operation within that module,
a \texttt{PoolQuery} encoding the routing mode and parallelism,
a \texttt{task\_group} for affinity-based scheduling, and
\texttt{task\_flags} that mark properties such as periodicity or
data ownership.
Derived task classes (e.g., \texttt{bdev::WriteTask}) add
operation-specific parameters---buffer pointers, offsets,
sizes---and implement \texttt{SerializeIn}/\texttt{SerializeOut}
methods for portable cross-architecture transfer.

The \texttt{PoolQuery} attached to each task controls three
dimensions of execution.
\emph{Routing mode} determines where the task executes: on a
specific GPU (\texttt{ToLocalGpu}), on the CPU
(\texttt{ToLocalCpu}), via hash-based load balancing
(\texttt{DirectHash}), or broadcast to all instances.
\emph{Parallelism} specifies the number of GPU threads the task
requires (default~1; e.g., 32 for warp-level block I/O).
The CPU runtime ignores this field; the GPU runtime uses it to
determine cross-warp scheduling, as described in
\S\ref{sec:orchestrator}.

%%%%============================================================
\subsection{Dynamic GPU Work Orchestration}
\label{sec:orchestrator}
%%%%============================================================

In a shared GPU runtime, multiple clients may simultaneously submit
tasks with widely varying parallelism requirements.  Furthermore,
machines experience periods of high and low demand; a persistently
executing GPU kernel can waste energy during idle periods.
The work orchestrator addresses both challenges through hierarchical
load tracking, a multi-level task dispersion algorithm, and an
adaptive rescaling mechanism.

\subsubsection{Hierarchical Load Tracking}
\label{sec:load-tracking}

Let $W$ denote the total number of warps allocated to the GPU
runtime, and let $G = \lceil\sqrt{W}\rceil$ be the number of
\emph{warp groups}.  Each warp group contains $G$ warps, forming
a two-level hierarchy indexed by a \emph{group ID}
$g = \lfloor\mathit{warp\_id} / G\rfloor$ and a \emph{minor ID}
$m = \mathit{warp\_id} \bmod G$.

Load is tracked in two shared-memory tables residing in GPU global
memory:

\begin{itemize}[nosep]
  \item \texttt{warp\_group\_load\_[$G$]}:
    an array of atomic 32-bit counters, one per group, recording
    the aggregate number of active tasks in the group.
  \item \texttt{warp\_load\_[$G$][$G$]}:
    a two-dimensional array recording the number of active tasks
    in each individual warp, indexed by
    $[\mathit{group\_id}][\mathit{minor\_id}]$.
\end{itemize}

\noindent
When a warp begins executing a task whose \texttt{Future} range
fits within a single warp (i.e., range width $\leq 32$), it
atomically increments both its group-level and warp-level counters.
When the task completes, it decrements both.
Because each warp updates only its own entry in the warp-level
table, contention is limited to the group-level counters, which
are updated at most once per task arrival and departure.
The $\sqrt{W}$ grouping ensures that the number of groups remains
small (e.g., 8~groups for 64~warps), keeping the cost of scanning
group loads constant.

\subsubsection{Multi-Level Work Assignment}
\label{sec:work-assign}

When a new task arrives at the runtime, it is first placed on a
per-warp queue using round-robin selection from a random offset.
The receiving warp then inspects the task's parallelism requirement
$P$ (from \texttt{PoolQuery.parallelism}) and dispatches it through
one of three levels of the scheduling hierarchy.
Algorithm~\ref{alg:dispatch} summarizes the procedure.

\begin{algorithm}[t]
\caption{Multi-level task dispatch}
\label{alg:dispatch}
\begin{algorithmic}[1]
\Require Task $t$ with parallelism $P$, arriving at warp group $g$
\Ensure $t$ is assigned to the appropriate warps
\State $G \gets \lceil\sqrt{W}\rceil$ \Comment{warps per group}
\If{$P \leq 32$} \Comment{\textbf{Case 1: Warp-level}}
    \State Allocate coroutine frame on local warp stack
    \State Execute $t$ directly on this warp
    \State \texttt{warp\_group\_load\_[$g$]}\texttt{++}
    \State \texttt{warp\_load\_[$g$][$m$]}\texttt{++}
\ElsIf{$P \leq 32 \cdot G$} \Comment{\textbf{Case 2: Group-level}}
    \State $N_w \gets \lceil P / 32 \rceil$
    \State Select $N_w$ least-loaded warps within group $g$
    \For{each selected warp $w_i$ with rank $i$}
        \State Create \texttt{Future} copy with range $[32i,\; 32(i+1))$
        \State Enqueue copy on $w_i$'s task queue
    \EndFor
\Else \Comment{\textbf{Case 3: Multi-group}}
    \State $N_g \gets \lceil P / (32 \cdot G) \rceil$
    \State Select $N_g$ least-loaded groups, preferring groups
           nearest to $g$
    \For{each selected group $g_j$ with rank $j$}
        \State Create \texttt{Future} copy with range
               $[32 G j,\; 32 G (j+1))$
        \State Enqueue copy on $g_j$'s group queue
        \State Group $g_j$ recurses to \textbf{Case~2}
    \EndFor
\EndIf
\end{algorithmic}
\end{algorithm}

\paragraph{Case 1: Warp-level tasks ($P \leq 32$).}
The common case for fine-grained I/O operations.  The warp allocates
a coroutine frame from its stack and executes the task directly.
No cross-warp coordination is needed.

\paragraph{Case 2: Group-level tasks ($32 < P \leq 32G$).}
The task requires multiple warps but fits within a single group.
The scheduler computes $N_w = \lceil P/32 \rceil$ and selects the
$N_w$ least-loaded warps within the local group by scanning
\texttt{warp\_load\_[$g$][$\cdot$]}.
It then creates $N_w$ shallow copies of the task's \texttt{Future},
each annotated with a disjoint \emph{range} indicating which
subset of the work the warp is responsible for (e.g., which
output elements to compute).  All copies share the same
underlying \texttt{FutureShm}, avoiding data duplication.
The module accesses \texttt{GetTaskWarpOffset()} to determine its
assigned partition of the computation.

\paragraph{Case 3: Multi-group tasks ($P > 32G$).}
The task exceeds the capacity of a single group.  The scheduler
computes $N_g = \lceil P/(32G) \rceil$ and selects the $N_g$
least-loaded groups by scanning \texttt{warp\_group\_load\_},
starting from the local group and spiraling outward to preserve
locality.  Each selected group receives a \texttt{Future} copy
whose range spans $32G$ threads; the group then internally
dispatches via Case~2.

\paragraph{Cross-warp completion.}
When a task is fragmented across multiple warps, each warp
independently executes its partition and, upon finishing,
atomically increments a \texttt{completion\_counter} stored in
the shared \texttt{FutureShm}.  The warp whose increment reaches
$P$ is designated the \emph{completer}: it copies any return
values back to the client (for client-originated tasks) or signals
the parent coroutine to resume (for runtime-internal sub-tasks).
This mechanism avoids any barrier synchronization; warps complete
asynchronously and the last one to finish performs cleanup.

\subsubsection{Adaptive Runtime Rescaling}
\label{sec:rescaling}

To avoid wasting GPU resources during idle periods, the orchestrator
periodically evaluates aggregate load and adjusts the number of
active warps.  The rescaling algorithm operates on a configurable
interval $\Delta t$ (default: 100\,ms) and proceeds as follows:

\begin{enumerate}[nosep]
  \item Compute $L = \sum_{g} \texttt{warp\_group\_load\_}[g]$,
    the total number of active tasks across all groups.
  \item If $L = 0$ for $k$ consecutive intervals (default $k=10$),
    the runtime \emph{parks} idle warp groups by setting their
    polling flags to a low-frequency spin, reducing energy
    consumption.
  \item If $L > \alpha \cdot W$ (high-watermark, default
    $\alpha = 0.8$), all parked groups are reactivated.
  \item If $L < \beta \cdot W_{\text{active}}$ (low-watermark,
    default $\beta = 0.2$), excess groups are parked, down to a
    minimum of one active group.
\end{enumerate}

\noindent
The asymmetric thresholds prevent oscillation: groups are activated
eagerly but deactivated conservatively.  Because parking only reduces
polling frequency (rather than terminating warps), reactivation is
instantaneous---no kernel relaunch is required.

%%%%============================================================
\subsection{AggroFS: A GPU-Resident Filesystem}
\label{sec:aggrofs}
%%%%============================================================

To demonstrate the utility of \sysname{}'s execution engine, we
build \emph{AggroFS}, a GPU-resident filesystem that manages metadata
and storage tiering entirely on the GPU\@.
The key insight is that storage tiering---placing data on the
fastest available device with sufficient capacity---requires
\emph{on-demand} knowledge of device occupancy and metadata state.
Prior work pre-allocates metadata on the CPU and communicates
placement decisions across the PCIe bus; in workloads where metadata
changes rapidly (e.g., checkpoint/restart, shuffle phases), this
communication becomes a bottleneck.  By co-locating metadata with the
GPU compute that consumes it, AggroFS eliminates this round-trip
entirely.

\subsubsection{The Blob Store}
\label{sec:blobstore}

AggroFS adopts a \emph{blob store} API internally and layers a
POSIX-like filesystem interface on top.  The store is built on two
primitives: \emph{blobs} and \emph{tags}.

A \textbf{blob} is an uninterpreted byte sequence of arbitrary
size, divided into one or more fixed-size \emph{blocks} (default
1\,MB) that may reside on different storage tiers.  Each blob
carries a \emph{score} in $[0,1]$ that encodes its relative
importance for tiering decisions (higher scores prefer faster
tiers), along with timestamps for last modification and last read.

A \textbf{tag} is a named handle that gives meaning to one or more
blobs.  Tags can represent directories, filenames, or
application-defined groupings.  Each tag is identified by a
human-readable name and a system-assigned \texttt{TagId}.

Metadata is stored in three GPU-resident hash maps with
fine-grained, per-bucket cooperative locks:

\begin{itemize}[nosep]
  \item \texttt{tag\_name\_to\_id}: maps tag names to
    \texttt{TagId}s.
  \item \texttt{tag\_id\_to\_info}: maps \texttt{TagId}s to
    \texttt{TagInfo} records (name, total size, timestamps).
  \item \texttt{blob\_map}: maps composite keys
    \texttt{(tag\_id, blob\_name)} to \texttt{BlobInfo} records
    (block list, score, compression metadata).
\end{itemize}

\noindent
The maps use a fixed number of buckets configured at startup (default:
1\,M entries for blobs, 100\,K for tags).  Locks are acquired only
during metadata \emph{creation}; lookups and score updates proceed
lock-free using atomic compare-and-swap on the map's internal
linked-list pointers.

\subsubsection{Metadata Scalability and Eviction}
\label{sec:meta-evict}

GPU memory is a scarce resource; AggroFS must ensure that metadata
does not overwhelm device memory as the dataset grows.  We address
this through \emph{partitioned ownership} and \emph{LRU eviction}.

Each warp in the GPU runtime owns a fixed partition of the blob and
tag maps, determined by hashing the tag or blob key to a warp ID\@.
Ownership means that only the owning warp may create, modify, or
evict entries in its partition; other warps route metadata requests
to the owner via the intra-GPU task queues.
This partitioning eliminates global locks on the metadata maps:
concurrent creates on different keys land on different warps and
proceed without contention.

When a warp's partition reaches its capacity limit, the warp triggers
LRU eviction.  It scans its local partition for the entry with the
oldest \texttt{last\_read} timestamp, flushes any dirty blocks to
a lower storage tier (see \S\ref{sec:tiering}), and reclaims the
metadata slot.  Because eviction is warp-local, it does not block
other warps and can overlap with ongoing I/O on other partitions.

\subsubsection{Storage Tiering and Data Placement}
\label{sec:tiering}

Users register one or more block devices with AggroFS, each
annotated with a capacity limit, a persistence level
(\emph{volatile}, \emph{temporary}, or \emph{long-term}), and an
optional performance score.  Device types include GPU high-bandwidth
memory (HBM), pinned host memory, RAM disks, and NVMe SSDs.

\paragraph{Data placement.}
When a blob is written, AggroFS selects a target device using a
\emph{maximum-bandwidth, minimum-latency} policy.  For large
requests ($\geq 32$\,KB), devices are ranked by write bandwidth;
for small requests, by access latency.  Within each ranking, the
algorithm filters to devices with sufficient free capacity and
selects the highest-ranked device whose score does not exceed the
blob's score.  If no such device exists, the algorithm falls back
to lower-score tiers in decreasing-bandwidth order.

\paragraph{Blob rescoring and reorganization.}
Blob scores evolve over time.  AggroFS supports a
\texttt{ReorganizeBlob} operation that updates a blob's score and,
if the new score warrants a different tier, asynchronously migrates
the blob's blocks.  Migration proceeds block-by-block: each block
is read from its current device, written to the new target, and the
blob's block list is updated atomically.  The blob remains readable
throughout migration because the old blocks are freed only after
the new blocks are committed.

Because the GPU has direct, low-latency access to device occupancy
counters and blob metadata, tiering decisions can be made
\emph{inline} during the write path without a PCIe round-trip.
This is the central advantage of GPU-resident metadata: placement is
always computed against the most current view of storage state.

\subsubsection{Block Allocation}
\label{sec:block-alloc}

Each registered block device maintains a \emph{block allocator}
partitioned between GPU and CPU ownership.  The GPU partition serves
GPU-initiated allocations; the CPU partition serves CPU-initiated
ones.

Within each partition, the allocator combines two mechanisms.
A \emph{per-warp free-list cache} holds recently freed blocks,
organized by size category (256\,B through 1\,MB in six classes).
Allocations first consult the local cache; only on a miss do they
fall through to a global \emph{bump allocator} backed by an atomic
heap pointer.  The bump allocator returns a contiguous range from
the device's address space, which is carved into blocks of the
requested size category.

When the GPU partition is exhausted, the allocator sends a
\texttt{ToLocalCpu} task requesting a bulk transfer of free blocks
from the CPU's pool.  Conversely, an exhausted CPU partition
requests blocks from the GPU\@.  This cross-partition rebalancing
is coarse-grained (transferring many blocks at once) to amortize
the PCIe latency of the request.

\subsubsection{A Filesystem on the Blob Store}
\label{sec:fs-layer}

AggroFS exposes a POSIX-like interface by mapping filesystem
abstractions onto the blob store.  Tags represent absolute
paths (directories and filenames); files are decomposed into
a sequence of blobs, each of configurable size (default 1\,MB),
keyed by \texttt{(tag\_id, chunk\_index)}.

To amortize the cost of small writes, AggroFS supports
\emph{blob preallocation}: when a file is opened for writing,
a configurable number of blob slots are reserved in the metadata
map and their blocks are pre-allocated on the target device.
Subsequent writes fill pre-allocated blobs without acquiring any
locks, falling back to on-demand allocation only when the
preallocated pool is exhausted.

Reads and writes are fully asynchronous.  A GPU application
issues \texttt{PutBlob} or \texttt{GetBlob} tasks and
receives a \texttt{Future} that can be \texttt{co\_await}ed
in a warp-level coroutine.  This enables applications to
overlap I/O with computation at fine granularity---for example,
a training pipeline can prefetch the next mini-batch's data
while the current batch is being processed.

\subsubsection{Metadata Consistency and Persistence}
\label{sec:consistency}

AggroFS maintains metadata consistency through \emph{asynchronous
background flushing}.  Each warp runs a low-priority periodic task
that scans its metadata partition for dirty entries---blobs or tags
whose in-GPU state has diverged from the CPU-side persistent copy.
Dirty entries are serialized and sent to the CPU runtime via
\texttt{ToLocalCpu} tasks, where they are written to durable
storage.

The flushing protocol proceeds in three steps:
(1)~the warp snapshots the dirty entry's metadata and marks it
as \emph{flush-pending};
(2)~the CPU runtime writes the snapshot to the persistent store
and acknowledges completion;
(3)~the warp clears the dirty and flush-pending flags.
If the entry is modified while flush-pending, it is re-marked
dirty and will be flushed again in the next cycle.  This design
ensures that the persistent copy is always a consistent (though
possibly stale) prefix of the GPU's state.

Flush frequency is configurable per-warp, allowing operators to
trade durability for throughput.  For volatile tiers (e.g., HBM
scratch space), flushing can be disabled entirely.

\subsubsection{Live Module Upgrade}
\label{sec:live-upgrade}

\sysname{} supports upgrading module code (e.g., deploying a new
version of the AggroFS filesystem logic) without losing in-flight
data.  The upgrade protocol proceeds as follows:

\begin{enumerate}[nosep]
  \item \textbf{Quiesce.}  The orchestrator stops accepting new
    tasks for the target module and drains all in-flight tasks to
    completion.
  \item \textbf{Flush.}  All dirty metadata and data in volatile
    tiers are flushed to the CPU runtime and persisted to durable
    storage, following the protocol of \S\ref{sec:consistency}.
  \item \textbf{Reload.}  The runtime unloads the old module
    shared library and loads the new version.  On the GPU, this
    updates the inline function-pointer table in the container
    object.  On the CPU, the virtual dispatch table is replaced by
    the new library's symbols.
  \item \textbf{Reinitialize.}  The new module's \texttt{Create}
    callback is invoked with a flag indicating an upgrade rather
    than a fresh start.  The module reloads its persistent
    metadata from durable storage and reconstructs any in-memory
    indices.
  \item \textbf{Resume.}  The orchestrator re-enables task
    acceptance for the module.
\end{enumerate}

\noindent
The runtime is fully restarted during step~3 to ensure a clean GPU
state; however, because all volatile data was flushed in step~2, no
user data is lost.  The total upgrade latency is dominated by the
flush phase, which is bounded by the volume of dirty metadata and
the throughput of the durable storage tier.
